% !TEX root = ../test_multifile.tex
\section{Chapter Three: Results and Discussion}

In this chapter we present the experimental results and discuss their
implications. The data collected from all three benchmark datasets are
summarized below, followed by a detailed comparison with existing methods
from the literature.

The synthetic dataset produced exact recovery of all frequency components
up to the Nyquist limit, confirming the correctness of our implementation.
Peak signal-to-noise ratios exceeded 120 dB for all test cases, which is
consistent with double-precision floating-point arithmetic.

For the audio dataset, our method identified dominant frequencies within
0.1 Hz of the ground truth values obtained from a professional spectrum
analyzer. The total processing time for a one-hour recording sampled at
44100 Hz was 3.2 seconds on commodity hardware, demonstrating the
practical viability of the approach.

The financial time series revealed interesting periodic patterns at
multiple scales. Weekly and monthly cycles were clearly visible in the
spectral decomposition. Longer-period oscillations corresponding to
quarterly earnings reports and annual seasonal effects were also detected,
though with lower statistical significance.

We conclude that the method is both accurate and efficient for a wide range
of signal types. Future work will investigate extensions to higher dimensions
and adaptive windowing strategies for non-stationary signals.
